\documentclass[10pt,a4paper]{article}
\usepackage[margin=0.40in,top=0.38in,bottom=0.42in]{geometry}
\usepackage[T1]{fontenc}
\usepackage{newtxtext}
\usepackage{enumitem}
\usepackage{hyperref}
\usepackage{microtype}
\hypersetup{colorlinks=true,urlcolor=blue,linkcolor=blue,pdftitle={Resume - Sehyun Choi},pdfauthor={Sehyun Choi}}
\pagestyle{plain}
\setlength{\parindent}{0pt}
\setlength{\parskip}{0pt}
\setlength{\emergencystretch}{1.2em}
\setlist[itemize]{leftmargin=1.15em,label=\textbullet,itemsep=0.5pt,topsep=0.5pt,parsep=0pt,partopsep=0pt}
\newcommand{\sectiontitle}[1]{%
  \par\vspace{4pt}%
  \noindent{\bfseries\fontsize{10.3}{12}\selectfont #1}\par
  \vspace{1.5pt}\hrule\vspace{2pt}%
}
\newcommand{\job}[4]{\noindent\textbf{#1}\hfill\textbf{#2}\\[0.5pt]\textit{#3}\hfill #4\par\vspace{1pt}}
\newcommand{\eduline}[4]{\noindent\textbf{#1}\hfill\textbf{#2}\\[0.5pt]\textit{#3}\hfill #4\par\vspace{1pt}}

\begin{document}
\fontsize{9.35}{10.65}\selectfont
\begin{center}
  {\fontsize{20}{22}\selectfont\bfseries Sehyun CHOI}\\[1pt]
  Address: Seoul South Korea \\[1pt]
  Email: \href{mailto:thyconds@gmail.com}{thyconds@gmail.com} | Telephone: +82 01088194520 | Website: \href{https://syncdoth.github.io}{syncdoth.github.io}
\end{center}
\vspace{-5pt}

\sectiontitle{EDUCATION}
\eduline{The Hong Kong University of Science and Technology (HKUST)}{Hong Kong}{Master of Philosophy (MPhil) in Computer Science, supervised by Prof. Yangqiu Song}{09/2022 - 09/2024}
\begin{itemize}
  \item \textbf{CGA}: 4.12/4.30 (Asian Future Leaders Scholarship Awardee)
\end{itemize}
\textit{Bachelor of Engineering (BEng) in Computer Science, Minor in Bioengineering}\hfill 09/2017 - 06/2022
\begin{itemize}
  \item \textbf{GGA}: 4.01/4.30 (First Class Honours; Dean's List and Academic Achievement Medal Awardee)
  \item \textbf{Relevant Coursework}: Machine Learning for NLP, Deep Learning in Computer Vision, Big Data Management, Honours Software Engineering, Honours Algorithms, Honours Object Oriented Programming and Data Structures, Linear Algebra
\end{itemize}
\eduline{Handong International School}{South Korea}{Advanced Placement (AP) by College Board and ACT}{02/2011 - 02/2016}
\begin{itemize}
  \item \textbf{Subject}: Physics I (5/5), Physics II (5/5), Calculus BC (5/5) | ACT (34/36)
\end{itemize}

\sectiontitle{PROFESSIONAL \& RESEARCH EXPERIENCE}
\job{TwelveLabs}{South Korea}{ML Research Scientist}{10/2024 - Present}
\begin{itemize}
  \item Implemented video-native frame selection and patch tokenization technology, effectively reducing the maximum token budget by 50\% while increasing the average benchmark score by 5\%p.
  \item Developed an effective research consolidation protocol with Multi-teacher On-Policy Distillation, enabling parallel research work and integration into a single release weight.
  \item Improved video-LLM training efficiency 2x by improving distributed training communications, video token load balancing, and incorporating \& writing cutting-edge GPU kernels.
  \item Improved Video-LLM on-policy RL efficiency 4x for video inputs with overlapped communication, node-aware rollout routing, and efficient multimodal cache scheduling and communications.
\end{itemize}
\job{NucleusAI}{U.S., Remote}{Research Intern}{09/2023 - 07/2024}
\begin{itemize}
  \item Implemented from scratch a high-inference-speed language model based on the Retentive Network architecture and pretrained it with scalable training techniques such as zero-redundant data parallelism, pipeline parallelism, and tensor parallelism across multi-GPU and multi-node systems.
\end{itemize}
\job{Naver Corporation}{South Korea}{Machine Learning Research Intern (Papago)}{12/2021 - 02/2022}
\begin{itemize}
  \item Experimented with Deep Active Learning approaches for querying the most critical data points from unlabeled datasets in Computer Vision and Natural Language Understanding for Machine Translation (MT) tasks.
  \item Improved fine-tuning efficiency by selecting informative data based on corpus statistics, pretrained MT Transformer output uncertainty, and the model's sentence encoder representations.
  \item Identified uncertainty, representativeness, and diversity as essential aspects of Active Learning data selection, and devised metrics to quantify them.
\end{itemize}
\job{KAIST Artificial Intelligence (AI) Lab}{South Korea}{Summer Research Intern}{07/2021 - 08/2021}
\begin{itemize}
  \item Supported a language model detoxification project by researching eXplainable AI methods to detect potentially offensive concepts captured in language models and ablate them.
  \item Proposed a novel framework for detecting offensive concepts in language models from non-offensive sentences.
\end{itemize}
% \job{HKUST Undergraduate Research Opportunity Program (UROP)}{Hong Kong}{Student Researcher (\href{https://github.com/HKUST-KnowComp/CSKB-Population}{project})}{02/2021 - 12/2021}
% \begin{itemize}
%   \item Developed novel methods using Graph Neural Networks (GNNs) and pretrained language models (PTLMs) on Commonsense Knowledge Graphs (CKGs), achieving significant improvement over previous baselines in the CKG Population task.
%   \item Performed hyperparameter optimization and ablation studies over network components such as PTLMs and GNN designs to find the best-performing model and quantify each component's contribution.
% \end{itemize}
\job{Skelter Labs}{South Korea}{Machine Learning Engineer Intern}{12/2019 - 07/2020}
\begin{itemize}
  \item Conducted experiments to improve deep learning-based speech synthesis models, including introducing new network modules, using data augmentation, and fine-tuning on different datasets.
  \item Used data pipelines, microservices, and cloud services such as Spark, Kubernetes, and GCP to run large-scale training jobs.
  \item Participated in collaborative software development using Git and code review systems.
\end{itemize}

\sectiontitle{PUBLICATIONS}
\textbf{KCTS: Knowledge-Constrained Tree Search Decoding with Token-Level Hallucination Detection}\\[-1.5pt]
\textit{EMNLP 2023 Main Conference} | \href{https://arxiv.org/abs/2310.09044}{arxiv.org/abs/2310.09044}\par
\textbf{AbsPyramid: Benchmarking the Abstraction Ability of Language Models with a Unified Entailment Graph}\\[-1.5pt]
\textit{Arxiv Preprint} | \href{https://arxiv.org/abs/2311.09174}{arxiv.org/abs/2311.09174}\par
\textbf{CKBP v2: An Expert-Annotated Evaluation Set for Commonsense Knowledge Base Population}\\[-1.5pt]
\textit{Arxiv Preprint} | \href{https://arxiv.org/abs/2304.10392}{arxiv.org/abs/2304.10392}\par
\textbf{Benchmarking Commonsense Knowledge Base Population with an Effective Evaluation Dataset}\\[-1.5pt]
\textit{EMNLP 2021 Main Conference} | DOI: \href{https://doi.org/10.18653/v1/2021.emnlp-main.705}{10.18653/v1/2021.emnlp-main.705}

\sectiontitle{AWARD \& CERTIFICATE}
\job{Naver Clova AI Rush 2021 \& 2022}{South Korea}{2nd place in Unknown Document Classification (2022) and Smart Grammar Editor (2021)}{2021, 2022}
\begin{itemize}
  \item Trained Korean pretrained language models with model calibration and representation-distance-based out-of-distribution detection to identify unknown document distributions while maintaining in-domain classification performance.
  \item Improved the benchmark GLEU score by 33\% over the vanilla Transformer baseline for Korean Grammatical Error Correction (GEC) by implementing pseudo-data creation with back-translation, two-stage pretraining, and data mixing during fine-tuning.
\end{itemize}
% \job{EY Next Wave Data Science Competition}{Hong Kong}{Country Finalist - Hong Kong (\href{https://github.com/syncdoth/EY_DataWave_Challenge}{project})}{04/2019 - 05/2019}
% \begin{itemize}
%   \item Performed feature engineering and cleaned raw data, then developed a Recurrent Neural Network (LSTM) to predict city traveler behavior within a specific time window.
% \end{itemize}

% \sectiontitle{EXTRACURRICULAR ACTIVITY}
% \job{HKUST KSA Machine Learning Study Club}{Hong Kong}{Founding Member \& President (\href{https://github.com/syncdoth/ML_STUDY_2020}{club})}{10/2020 - 12/2021}
% \begin{itemize}
%   \item Organized a machine learning study club of more than 20 students and taught a self-developed curriculum covering Python, linear algebra, deep learning with the TensorFlow API, and real-world projects using state-of-the-art models.
% \end{itemize}
% \job{J.P. Morgan - Code for Good 2020}{Hong Kong}{Participant}{10/2020}
% \begin{itemize}
%   \item Designed a NoSQL schema for Firestore and React front-end components to visualize student submissions in an administrative website for centralized communication and progress reporting.
% \end{itemize}

\sectiontitle{SKILL \& INTEREST}
\textbf{Language}: Korean (Native) | English (Fluent) | Python\\
\textbf{Interest}: NLP | LLM | Controllable Generation | Knowledge Grounding \& Reasoning | AI Safety | XAI
\end{document}
